% figures/pd-buildid.tex -- \pdbuildid: which build this PDF is.
%
%   \pdbuildid  ->  "dfce4990d · open-proof · 9 November 2026"
%
% The short commit the Book was built from, the font profile, and the date of
% SOURCE_DATE_EPOCH. It is meant for the colophon and the how-to-cite line,
% where a hand-written date ("5 November 2026" in textbook.json) goes stale
% the day after it is typed; placing it there is the front matter's call.
%
% Where the parts come from:
%   \pdbuildcommit  scripts/generate-mega-whitepaper.mjs, from Git (HEAD, with
%                   "+modified" when a Book input has uncommitted changes), or
%                   from SOURCE-MANIFEST.json in the source archive, or from
%                   PD_BOOK_COMMIT when it is set;
%   \pdbuilddate    the same generator, from SOURCE_DATE_EPOCH, which
%                   scripts/build-whitepapers.sh computes per build and the
%                   source archive pins;
%   \pdbuildfontprofile  scripts/prepare-book-fonts.py's \pdbookfontprofile
%                   (open-proof or suisse), read when \pdbuildid is used.
% The generator writes the first two into mega-volume-build-id.tex. Outside a
% Book build (a figure compiled alone) that file may be absent, and the parts
% say "unversioned" and "undated" rather than stopping the compile.
\IfFileExists{../.cache/whitepaper-build/coordination-papers-mega-volume/mega-volume-build-id.tex}
  {\input{../.cache/whitepaper-build/coordination-papers-mega-volume/mega-volume-build-id.tex}}{}
\providecommand\pdbuildcommit{unversioned}
\providecommand\pdbuildepoch{}
\providecommand\pdbuilddate{undated}
\newcommand\pdbuildfontprofile{\ifdefined\pdbookfontprofile\pdbookfontprofile\else unknown\fi}
\newcommand\pdbuildid{\pdbuildcommit\ \textperiodcentered\ \pdbuildfontprofile\ \textperiodcentered\ \pdbuilddate}
